\section{Model and funding frontiers}\label{sec:model}
An institution fixes a finite set $N$ of actual owners and $d$ resource types.
Owner $i$ has an endowment $\bar b_i\in\R_+^d$ and may report any
$b_i\in[0,\bar b_i]$ coordinatewise. Zero endowment coordinates encode
ineligibility. These owners, coordinates, boxes and the public technology
remain fixed throughout every comparison. Write $b_C=\sum_{i\in C}b_i$,
$B=b_N$, and $\mathcal B=\prod_{i\in N}[0,\bar b_i]$.
Different institutions may have different, arbitrarily large finite owner and
resource counts.

The technology $F:\R_+^d\to\R_+$ is finite-valued, monotone, concave and
homogeneous of degree one, with $F(0)=0$. These properties imply
superadditivity:
\begin{equation}\label{eq:superadditive}
 F(x+y)=2F((x+y)/2)\ge F(x)+F(y).
\end{equation}
No closed-orthant continuity assumption is imposed on the general dimension
class. In the sparse-production sections we specialize to a finite packing
program
\begin{equation}\label{eq:packing}
 F(c)=\max\{\pi\cdot x: Ax\le c,\ x\ge0\},\qquad A\ge0,\quad\pi\ge0.
\end{equation}
Every positive-profit column consumes a positive amount of at least one
resource. Zero-profit columns may be discarded. The optimum is finite and
attained at finite capacities. Its \emph{activity arity} is the maximum number
of positive entries in a positive-profit column. Arity three permits many
resource rows; it does not mean three resources in total.

A rule is an arbitrary map $y:\mathcal B\to\R_+^N$.
It need not be continuous, anonymous, marginal, or linear. All requirements
below concern the same rule at every real report profile. For
$y_C(b)=\sum_{i\in C}y_i(b)$, \emph{coalition coverage} is
\begin{equation}\label{eq:coverage}
 y_C(b)\ge F(b_C)\qquad(C\subseteq N,\ b\in\mathcal B).
\end{equation}
Coverage is required for \emph{every} coalition, without a size restriction.
A total-funding factor $\kappa$ means
\begin{equation}\label{eq:funding}
 \sum_{i\in N}y_i(b)\le\kappa F(B)\qquad(b\in\mathcal B).
\end{equation}
Thus the external subsidy is at most $(\kappa-1)F(B)$; $\kappa$ itself is a
factor on total payments. When $F(B)=0$, nonnegativity and
\eqref{eq:funding} force every payment to vanish. At a positive-production
profile, an owner reporting zero is allowed to receive a positive payment.
There is no deletion-marginal floor or separate dummy-zero axiom.

\begin{definition}[Coalition withholding contracts]\label{def:contracts}
Fix an integer $s\ge1$. For every $C\subseteq N$ with $|C|\le s$, let
$\delta_i\in[0,b_i]$ for $i\in C$ and $\delta_i=0$ outside $C$.
\emph{Transferable coalition withholding monotonicity} (TCWM) requires
\begin{equation}\label{eq:tcwm}
 y_C(b)-y_C(b-\delta)\ge0.
\end{equation}
\emph{Full outside-production protection} (full OUT) requires
\begin{equation}\label{eq:out}
 y_C(b)-y_C(b-\delta)\ge F(\delta_C).
\end{equation}
Protected coalitions may contain unchanged members: $\delta_i=0$ is allowed
for any member of $C$. For unrestricted coalition protection, these
inequalities are required for every $C\subseteq N$.
\end{definition}

Both contracts permit side payments within the protected coalition. TCWM
uses payment-only utility; full OUT additionally allows the coalition to
produce with its pooled withheld resources. Neither is a model of arbitrary
misreports, private opportunity costs, or an endogenous change in the owner
roster. Full OUT implies TCWM. A lower bound for TCWM therefore also applies
to full OUT, whereas a TCWM construction need not satisfy full OUT.

For an institution $I$, let $k_s(I)$ and $k_s^{\mathrm{OUT}}(I)$ be the
infima of its feasible factors under these respective contracts and
\eqref{eq:coverage}. Define
\begin{equation}\label{eq:frontiers}
 K_r(s)=\sup_{I:\,\text{arity}(I)\le r} k_s(I),\qquad
 K_r^{\mathrm{OUT}}(s)=\sup_{I:\,\text{arity}(I)\le r}
 k_s^{\mathrm{OUT}}(I).
\end{equation}
The supremum ranges over all finite packing institutions, including their
actual eligibility and endowment boxes. A uniform upper constructs one rule
for each entire institution; a lower constructs an institution on which
\emph{every} rule has an expensive feasible report. The infimum convention
does not assume optimal rules are attained. We use $D_d$ and
$D_d^{\mathrm{OUT}}$ for the corresponding worst-instance factors for general
$d$-resource technologies when \emph{every} coalition receives protection.

\begin{remark}[Static feasibility versus cross-report incentives]\label{rem:static}
For a packing technology, at each fixed profile an optimal nonnegative dual
price $w$ gives payments $w\cdot b_i$. Dual feasibility implies
$w\cdot b_C\ge F(b_C)$, and their sum is $F(B)$.
Thus each individual report profile admits an efficient coalition cover.
The funding problem concerns selecting payments consistently across profiles.
Neither the static core nor an integrality-gap calculation by itself proves
our cross-report lower bounds.
\end{remark}

\section{Sparse production: a uniform coalition-size upper bound}\label{sec:upper}
We first state the upper bound with which the new lower must be compared.
The binomial semivalue and bounded-rank support inequalities are established
ingredients \citep{FeigeEtAl2014,EzraEtAl2018,Casajus2011}. Their use here retains
actual mixed owners, all-coalition coverage and the stronger full-OUT contract.
The proof is included to expose the precise comparison class.

\begin{theorem}[Sparse-production guarantee]\label{thm:sparse-upper}
For every finite packing institution of activity arity at most $r\ge1$ and
every integer $s\ge1$, there is an all-profile rule satisfying coalition
coverage and full OUT for every coalition of size at most $s$, with
\begin{equation}\label{eq:sparse-upper}
 K_r(s)\le K_r^{\mathrm{OUT}}(s)\le r s^{r-1}.
\end{equation}
For $p=1/s$, one such rule is
\begin{equation}\label{eq:semivalue}
 m_i(p;b)=\mathbb E_{R\subseteq N\setminus\{i\}}
       [F(b_i+b_R)-F(b_R)],\qquad y_i(b)=p^{1-r}m_i(p;b),
\end{equation}
where other owners are independently included in $R$ with probability $p$.
For rational data the rule is a finite rational computation. No
polynomial-time implementation or incentive guarantee for a sampled estimator
is asserted.
\end{theorem}

\begin{lemma}[Owner support and thinning]\label{lem:owner-support}
Fix a report profile and let $v(T)=F(b_T)$. For $i\in R$, set
$U_i(R)=v(R)-v(R\setminus\{i\})$. Then
\begin{align}
 \sum_{i\in T}U_i(R)&\ge v(T) &&(T\subseteq R),\label{eq:dual-cover}\\
 \sum_{i\in R}U_i(R)&\le r v(R),\label{eq:rank-loss}\\
 \mathbb E v(S\cap R_p)&\ge p^r v(S),\label{eq:thinning}
\end{align}
where $R_p$ independently retains every owner with probability $p$.
\end{lemma}
\begin{proof}
An optimal nonnegative dual price $w$ at $b_R$ satisfies
$v(R)=w\cdot b_R$ and $v(R\setminus\{i\})\le w\cdot b_{R\setminus\{i\}}$.
Thus $U_i(R)\ge w\cdot b_i$, and dual feasibility proves
\eqref{eq:dual-cover}.

For a fixed $S$, choose an optimal production plan $x$. For every used
activity $e$ and each required row $a$, assign that row to an actual supplier
$i\in S$ with probability $b_{ia}/(b_S)_a$, independently across required
rows. The denominator is positive whenever $x_e>0$ and $A_{ae}>0$.
Each assignment gives a nonnegative value piece, of weight
$\pi_e x_e$ times its assignment probability, supported on at most $r$ owners.
Let $h(T)$ be the total weight of pieces whose suppliers all belong to
$T\subseteq S$. For activity $e$, its retained fraction is
\[
 \prod_{a:A_{ae}>0}\frac{(b_T)_a}{(b_S)_a},
\]
which is at most the fraction in each individual required row. The retained
plan is therefore feasible for $b_T$. Consequently $h(T)\le v(T)$ and
$h(S)=v(S)$. Applying this at $S=R$ gives
\[
 U_i(R)\le h(R)-h(R\setminus\{i\});
\]
summing counts each value piece at most $r$ times and proves
\eqref{eq:rank-loss}. Under independent thinning, a piece supported on $k\le r$
owners survives with probability $p^k\ge p^r$. Hence
$\mathbb E v(S\cap R_p)\ge\mathbb E h(S\cap R_p)\ge p^r v(S)$.
This decomposes a production plan, not the strategic owners.
\end{proof}

\begin{proof}[Proof of Theorem~\ref{thm:sparse-upper}]
Let $R_p$ now thin all owners. Conditioning on whether $i\in R_p$ gives
\[
 p\sum_{i\in S}m_i(p;b)
 =\mathbb E\sum_{i\in S\cap R_p}U_i(R_p)
 \ge\mathbb E v(S\cap R_p)\ge p^r v(S).
\]
After multiplication by $p^{-r}$, this is coverage of every $S$, including
$|S|>s$. Nonnegativity follows from monotonicity. Also
\[
 p\sum_i m_i(p;b)\le r\mathbb E F(b_{R_p})
 \le rF(pB)=rpF(B),
\]
where the second inequality is concavity. This proves the funding bound.

For $s\ge2$, fix a protected coalition $C$ and write $M_C=\sum_{i\in C}m_i$.
Expanding the finite expectation yields
\[
 M_C(b)=\sum_{A\subseteq N}c_C(A)F(b_A),
\]
with, for $0<|A|<n$,
\begin{equation}\label{eq:coefficient}
 c_C(A)=p^{|A|-1}(1-p)^{n-|A|-1}(|A\cap C|-p|C|).
\end{equation}
The endpoints are $c_C(\varnothing)=-|C|(1-p)^{n-1}$ and
$c_C(N)=|C|p^{n-1}$. Since $p|C|\le1$, every coefficient of a value affected
by $C$'s report is nonnegative. For a withdrawal $\delta$ supported on $C$,
superadditivity gives
$F(b_A)-F(b_A-\delta_A)\ge F(\delta_A)$.
For $A\cap C=\varnothing$, both this difference and $F(\delta_A)$ vanish.
Therefore
\[
 M_C(b)-M_C(b-\delta)\ge\sum_A c_C(A)F(\delta_A)=M_C(\delta).
\]
The profile $\delta$ lies in the same endowment boxes, with the same owners
and the same rule. Coverage at that profile gives
$p^{1-r}M_C(\delta)\ge F(\delta_C)$, proving full OUT.

At $s=1$, $p=1$ and $m_i=F(B)-F(B-b_i)$. The second term is fixed under
$i$'s withdrawal, so superadditivity proves full OUT directly; the first two
inequalities of Lemma~\ref{lem:owner-support} give coverage and factor $r$.
If no positive-profit activity exists, the zero rule suffices. Empty
coalitions and the empty owner set are immediate.
\end{proof}

\section{A square-root lower bound for arity three}\label{sec:lower}
The following lower applies to arbitrary rules under the full report-domain
contract. It does not impose a marginal formula, an efficient-price floor,
anonymity, or a zero payment to inactive owners.

\begin{theorem}[Projective cap exchange]\label{thm:projective}
For each integer $M\ge1$, there is a unit-profit packing technology of arity
three on $d=2^{M+1}-1$ resources, with $n=2M+1$ fixed fully eligible owners,
for which every coalition-covering rule satisfying TCWM for coalitions of
size at most $2M-1$ has funding factor at least
\begin{equation}\label{eq:projective-finite}
 \max\left\{\frac23 Q_M,\ 1+\frac{2(M-1)}{3M}Q_M\right\},
 \qquad Q_M=\prod_{m=2}^M\frac{2m}{2m-1}
             =\frac{4^M}{2\binom{2M}{M}}.
\end{equation}
Every owner's report domain is the full real unit cube. Consequently, for
every integer $s\ge1$,
\begin{equation}\label{eq:main-interval}
 \sqrt{s/6}\le K_3(s)\le K_3^{\mathrm{OUT}}(s)\le3s^2.
\end{equation}
\end{theorem}

\subsection{Technology, legal symmetry and productive profiles}
Let $V=\mathbb F_2^{M+1}$ have basis $v,e_1,\ldots,e_M$.
Resources are the nonzero vectors of $V$. For every unordered triple
$\{a,b,a+b\}$ of distinct nonzero vectors, include one activity consuming one
unit of each row and yielding one unit of output. Duplicate triples are
identified. All activities have exactly three inputs.

There are actual owners $p,h_1,\ell_1,\ldots,h_M,\ell_M$, each endowed with
the all-ones vector. Every report in $[0,1]^d$ is legal for every owner.
The product of owner permutations and invertible binary linear maps acts on
this institution: the latter permute resources and activities, while owner
permutations preserve the identical report domains.

\begin{lemma}[Finite averaging]\label{lem:average}
If a rule on this institution satisfies coverage, size-$s$ TCWM (or full OUT)
and factor $\kappa$, then its finite group average satisfies the same
constraints and factor, and is equivariant. A lower on total payment of the
average at a profile of production one is a lower on the original rule at
some profile in that profile's finite symmetry orbit.
\end{lemma}
\begin{proof}
For a group element $g$ with owner permutation $\sigma_g$, define the
conjugate rule by $y^g(b)=\sigma_g^{-1}y(gb)$. Resource and owner maps preserve
coordinate order, report feasibility, coalition values, coalition sizes,
and the production value of a withheld bundle. Each contract inequality
therefore holds for every conjugate. All constraints are linear in payment
values, so their finite average also satisfies them. A change of variable
in the group sum proves equivariance. Finally, its total payment at $b$ is
the average of original total payments at the legal profiles $gb$, which
all have production $F(B)$. At least one is at least that average. This is a
proof transformation of rules; no strategic comparison changes an owner.
\end{proof}

For $m=1,\ldots,M$, define profile $C_m$ as follows. Owner $p$ reports the
unit vector in row $v$; each $h_j$ reports the unit vector in row $e_j$;
$\ell_j$ reports the unit vector in row $v+e_j$ if $j\le m$, and zero
otherwise. Every unmentioned resource coordinate is zero. In particular,
\emph{dormant high owners still report $e_j$} when $j>m$.

\begin{lemma}[Exact support and value]\label{lem:star-support}
The only activities supported at $C_m$ are the $m$ triples
$\{v,e_j,v+e_j\}$ with $j\le m$, and $F(C_m)=1$.
\end{lemma}
\begin{proof}
In the possible row set
$\{v\}\cup\{e_j,v+e_j:1\le j\le M\}$, the sum of private rows from
 distinct indices has two nonzero $e$ coordinates and leaves the set. The
sum of the private pair at index $j$ is $v$, and adding $v$ exchanges those
two rows. Thus every supported triple is precisely a complete active branch.
Every branch consumes the reported unit of $v$, giving the upper bound one;
one branch attains it. Equivalently, the full packing dual prices $v$ and
all zero-capacity rows at one and all other rows at zero. Supported columns
contain $v$ and unsupported columns contain a zero-capacity row, so this is
a feasible dual of value one for the full technology.
\end{proof}

\subsection{A constant seed and a fixed-owner payment ladder}
Work with the averaged rule of Lemma~\ref{lem:average}. At $C_1$, the linear
map exchanging $v$ and $e_1$ and fixing $e_j$ for $j>1$, together with the
owner swap $(p,h_1)$, preserves the whole profile. So does the invertible map
$v\mapsto v+e_1$, fixing every $e_j$, together with the swap $(p,\ell_1)$.
These maps fix the dormant high reports. Hence the three productive owners
receive a common payment $a$. Their coalition produces one, and coverage
requires $3a\ge1$. This uses static coverage for a three-member coalition,
not a three-member incentive constraint.

At $C_m$, permuting active indices and independently replacing
$e_j$ by $e_j+v$ for active $j$ preserves the entire profile after the matching
owner swaps. These invertible maps fix $v$ and all inactive high rows.
They act transitively on the $2m$ active private owners. Let $q_m$ be their
common payment and put $T_m=2m q_m$. The seed is
\begin{equation}\label{eq:seed}
 q_1\ge\frac13,\qquad T_1\ge\frac23.
\end{equation}

The transition $C_{m-1}\to C_m$ restores only the actual owner $\ell_m$.
Protect the coalition consisting of $\ell_m$ and the $2m-2$ old active
private owners. Its cardinality is $2m-1$, and every old member's report
is unchanged. At the old profile its payment is at least $T_{m-1}$:
$y_{\ell_m}(C_{m-1})$ is nonnegative, and is \emph{not} assumed to be zero.
At the new profile the protected sum is $(2m-1)q_m$. Therefore TCWM gives
\begin{equation}\label{eq:ladder}
 (2m-1)q_m\ge2(m-1)q_{m-1},\qquad
 T_m\ge\frac{2m}{2m-1}T_{m-1}.
\end{equation}
The unchanged new high owner $h_m$ and the cap owner $p$ are not in this
protected coalition. No owner is split, cloned, added, or removed.
Iterating \eqref{eq:ladder} proves $T_M\ge2Q_M/3$. Since production is one,
nonnegativity and funding imply $\kappa\ge2Q_M/3$.

Let $z=y_p(C_M)$. Terminal branch coverage gives $z+2q_M\ge1$, and total
payment is at least $z+2M q_M$. Hence
\[
 \kappa\ge z+2M q_M\ge1+2(M-1)q_M
 \ge1+\frac{2(M-1)}{3M}Q_M.
\]
This remains valid even if $2q_M>1$; it does not set $z=1-2q_M$.
For an asymmetric original rule, Lemma~\ref{lem:average} locates the costly
profile somewhere in the orbit, not necessarily at the displayed $C_M$.
This proves the finite assertion of Theorem~\ref{thm:projective}.

\subsection{The uniform exponent and an optimized tail}
The elementary estimate
\begin{equation}\label{eq:central-bound}
 Q_M\ge\frac12\sqrt{3M+1}
\end{equation}
follows by induction. The base $M=1$ is equality; after squaring, the step
reduces to
$4(M+1)^2(3M+1)-(2M+1)^2(3M+4)=M\ge0$.
For $M=\lfloor(s+1)/2\rfloor$, all incentive coalitions used have size at
most $s$, and $M\ge s/2$. Thus
$K_3(s)\ge2Q_M/3\ge\sqrt{3M+1}/3\ge\sqrt{s/6}$.
Combining this with Theorem~\ref{thm:sparse-upper} proves
\eqref{eq:main-interval}, including $s=1$.

\begin{corollary}[Optimized ladder certificate]\label{cor:tail}
Set $B_1=1$. For $s\ge2$, let $h=\lceil s/2\rceil$ and
\begin{equation}\label{eq:Bs}
 B_s=1+\frac{2(s-1)Q_h}{3h}
            \left(\frac{s-1}{s}\right)^{s-h}.
\end{equation}
Then $K_3(s)\ge B_s$, and
\begin{equation}\label{eq:tail-asymptotic}
 \lim_{s\to\infty}\frac{B_s}{\sqrt s}
   =\frac{\sqrt{2\pi}}{3\sqrt e}.
\end{equation}
The witness may use $2s+1$ fully eligible owners and $2^{s+1}-1$ resource
rows. In particular $B_6=493/243>2$, with 13 owners and 127 rows.
The optimization is of this ladder certificate, not of the unrestricted
funding frontier.
\end{corollary}
\begin{proof}
Keep $s$ fixed and permit arbitrary ladder length $M$. At stage $m$ protect
$\ell_m$ and any $k_m=\min\{s-1,2m-2\}$ old active private owners. Whole
pairs are unnecessary; all selected old owners have payments $q_{m-1}$ and
$q_m$ at the respective profiles. TCWM and nonnegativity imply
\[
 (k_m+1)q_m\ge k_m q_{m-1},\qquad
 q_M\ge\tfrac13 P_{M,s},\qquad
 P_{M,s}=\prod_{m=2}^M\frac{k_m}{k_m+1}.
\]
Terminal coverage therefore gives
$\kappa\ge1+2(M-1)P_{M,s}/3$.
Up to $h$ the uncapped recurrence gives $P_{h,s}=Q_h/h$; afterwards every
factor is $\alpha=(s-1)/s$. For $M\ge h$ and $M>1$, the consecutive ratio
of $(M-1)P_{M,s}$ is $M\alpha/(M-1)$, which is greater than one for $M<s$,
equal to one for $M=s$, and less than one for $M>s$.
Before the cap the increment increases; its $M=1$ value is zero and is
handled separately. Thus the equal maxima occur at $M=s,s+1$, proving
\eqref{eq:Bs}. The central-binomial asymptotic
$Q_h\sim\sqrt{\pi h}/2$, together with $h\sim s/2$ and
$\alpha^{s-h}\to e^{-1/2}$, proves \eqref{eq:tail-asymptotic}.
Finally $Q_3=8/5$ gives $B_6=493/243$.
\end{proof}

\subsection{Exact certificate form and what computation verifies}
The displayed argument has nonnegative rational multipliers. For fixed $M,s$,
use the rows
\[
 3q_1\ge1,\qquad (k_m+1)q_m-k_mq_{m-1}\ge0,
 \qquad z+2q_M\ge1,\qquad \kappa-z-2M q_M\ge0.
\]
Give the two final rows weight one. Starting with residual coefficient
$2(M-1)$ at $q_M$, eliminate backwards: weight the row at $m$ by the
residual divided by $k_m+1$, then replace the residual by that weight times
$k_m$. Weight the seed by the final residual divided by three.
All coefficients of $q_m$ and $z$ cancel, leaving
$\kappa\ge1+2(M-1)P_{M,s}/3$. At $M=1$ only the final rows are needed.
A lifted certificate can expand averaging back into the original
owner/profile variables; it retains the nonnegativity rows for dormant
payments rather than replacing those payments by zero.

The accompanying standard-library checks verify small full projective
technologies, the required linear maps, exact finite coefficient cancellation,
and the optimized scalar tail. A selected original-variable lift additionally
checks 34,440 inequalities cancelling 26,040 payment coordinates. These are
reproducible finite certificates supporting the written argument. The universal
group averaging and all-size geometry are conventional proofs, not statements
proved merely by the finite tests.

\section{Interpretation, scope and remaining questions}\label{sec:discussion}
The two parameters answer different questions. The dimension frontier allows
arbitrary technologies on $d$ resources and protects every coalition. The
sparse frontier limits each activity to $r$ resources, allows $d$ to grow,
and protects coalitions of at most $s$ owners while retaining coverage for all
coalitions. A resource-dimension upper cannot be substituted for an
owner-independent bound depending only on $r,s$.

At arity three, the square-root-to-quadratic gap remains open. The lower uses
exponentially many resource rows. Pure ownership is possible, but requires
more owners than the fully eligible construction; Appendix~\ref{sec:arity}
keeps the counts separate. Neither the exponent nor its sharper constant
claims minimal witness size. The simple binomial upper evaluates finitely
many coalition production programs, potentially exponentially many. Efficient
approximate evaluation requires separate incentive and coverage analysis.

The proposed contribution is the quantitative link between fixed-owner
coalition withholding and funding, particularly the projective cap exchange
that supplies a constant seed without a dummy-zero assumption. Projective
geometry, positive hypergraph supports, semivalues, and coalition monotonicity
are established ingredients. A specialist must assess whether a prior theorem,
possibly under different terminology or normalization, already implies the
present bounds. The accompanying priority matrix records inspected versions
and unresolved original-source comparisons. Internal model reviews and exact
certificates are not a historical-novelty certificate or human peer review.

A conditional infrastructure interpretation would give each owner genuine
rights over compute inputs and firm power capacity, with physically feasible
partial withdrawals and a sponsor who benefits from the stronger contract.
Public telemetry does not establish those rights. The earlier bounded
feasibility check supplied no compatible held-out outcome panel or sponsor
surplus estimate. Moreover, a Gaussian reserve model is not automatically an
arity-three packing technology. We therefore make no measured pooling,
reliability, energy-saving or net-value claim here.

This manuscript concerns coordinated withholding. It does not settle the
separate owner-count/resource-dimension gap for individual incentives in the
companion truthful-production project. Optional floor parameters, structured
production refinements and application measurement studies are not separate
new claims of this draft. The immediate questions are correctness of the
fixed-institution reductions, priority of the quantitative bounds, and whether
these results form a useful unit of publication.
